\subsection{Intervals and the recursive sampler}\label{sec:algorithm}
\subsubsection{Dyadic centers and radii}

Choose
\begin{equation}\label{eq:precision}
P=b+4\lceil\log_2(D+1)\rceil+8,\qquad
\varepsilon=2^{-P}.
\end{equation}
Use the reference gain \(g=\max\{1,S\}\), computed exactly from the weights.
Set \(r_0=1\) and \(c_{0,j}=0\). Obtain a dyadic \(q^\ast\) with
\(|q^\ast-\tanh(gr_{\ell-1})|\le\varepsilon/16\), and set
\(r_\ell=q^\ast+5\varepsilon/2\). Then
\begin{equation}\label{eq:rounded}
\tanh(gr_{\ell-1})+2\varepsilon
\le r_\ell\le\tanh(gr_{\ell-1})+3\varepsilon.
\end{equation}
A certified approximation followed by rational rounding supplies
this choice without finding the exact floor of a transcendental
number.

For each row store
\begin{equation}\label{eq:rowdata}
\sigma=\sum_j|w_{\ell,ij}|,\qquad
a=b_{\ell,i}+\sum_jw_{\ell,ij}c_{\ell-1,j},\qquad
\rho=\sigma r_{\ell-1},
\end{equation}
together with prefix sums of the absolute integer weight numerators.
Define the mathematical endpoints, midpoint, and image radius by
\[
y_\pm=\tanh(a\pm\rho),\qquad
m=\frac{y_++y_-}{2},\qquad R=\frac{y_+-y_-}{2}.
\]
Choose a dyadic \(c_{\ell,i}\in[-1,1]\) with
\(|c_{\ell,i}-m|\le\varepsilon\). Clamping a rounded center to
\([-1,1]\) preserves its error bound.
The exact \(m,R\) are scalar functions of stored rationals.
Online refinement does not recompute the row sum.

\begin{lemma}[Interval invariant]\label{lem:interval}
For every context, \(|h_{\ell,i}-c_{\ell,i}|\le r_\ell\).
With \(t=\tanh a\), \(q=\tanh\rho\),
\begin{equation}\label{eq:midrad}
m=\frac{t(1-q^2)}{1-t^2q^2},\qquad
R=\frac{q(1-t^2)}{1-t^2q^2}\le q.
\end{equation}
\end{lemma}
\paragraph{Proof idea.}
Propagate the endpoints of each row interval through monotone tanh. The row
norm controls its width, while a small rounding allowance keeps the stored
center and radius valid at the next layer.

\begin{proof}
The preactivation differs from \(a\) by at most
\(\sum_j|w_{\ell,ij}|r_{\ell-1}=\rho\), so monotonicity puts the
output in \([y_-,y_+]\). The tanh addition formula proves
\eqref{eq:midrad}. Its denominator is positive, and
\(1-t^2\le1-t^2q^2\). Hence
\[
R+|m-c_{\ell,i}|
\le\tanh(\sigma r_{\ell-1})+\varepsilon\le r_\ell.
\]
Induction starts with \(|x_j|=1\).
\end{proof}

\subsubsection{Normalized signs and the top output}

The internal sampler \(S_{\ell,i}\) returns a sign of mean
\(v_{\ell,i}=(h_{\ell,i}-c_{\ell,i})/r_\ell\).
At layer zero it returns \(x_i\).
For a nonzero row, choose \(j\) with probability
\(|w_{\ell,ij}|/\sigma\), request a fresh \(S_{\ell-1,j}\), and
multiply by \(\sgn(w_{\ell,ij})\).
This source has mean
\[
z=\sum_j\frac{w_{\ell,ij}}{\sigma}v_{\ell-1,j},\qquad
b_{\ell,i}+\sum_jw_{\ell,ij}h_{\ell-1,j}=a+\rho z.
\]
The next two sections construct a centered factory of mean
\begin{equation}\label{eq:centered}
G=\frac{\tanh(a+\rho z)-m}{R}.
\end{equation}

Put \(\alpha=R/r_\ell\), \(\kappa=(m-c_{\ell,i})/r_\ell\).
The internal sampler chooses the centered factory with probability
\(\alpha\), and constant signs with probabilities
\begin{equation}\label{eq:internalmixture}
\PP(+1\hbox{ constant})=\frac{1-\alpha+\kappa}{2},\qquad
\PP(-1\hbox{ constant})=\frac{1-\alpha-\kappa}{2}.
\end{equation}
The interval invariant gives \(\alpha+|\kappa|\le1\).
The resulting mean is \(\alpha G+\kappa=v_{\ell,i}\).

At the top use the actual midpoint and image radius:
\begin{equation}\label{eq:topmixture}
\PP(\hbox{centered factory})=R,\quad
\PP(+1\hbox{ constant})=\frac{1-R+m}{2},\quad
\PP(-1\hbox{ constant})=\frac{1-R-m}{2}.
\end{equation}
These probabilities are valid since \(|m|+R\le1\).
The sign has mean \(h_{D,1}\); return \((1+Y)/2\).
A zero row uses its known scalar probability directly.

\subsection{A positive majority factory}\label{sec:majority}

For \(\rho>0\), let
\(G_\rho(z)=\tanh(\rho z)/\tanh\rho\).
The unit-norm construction only needs \(\rho\le2\), but the
representation also supports the large-norm extension.
If \(M_k(z)\) is the mean of majority on \(2k+1\) independent signs
of mean \(z\), differentiating its binomial sum gives
\begin{equation}\label{eq:majderivative}
M_k'(z)=C_k(1-z^2)^k,\quad
C_k=\frac{(2k+1)\binom{2k}{k}}{4^k}\ge1,\quad
M_k(0)=0,\quad M_k(1)=1.
\end{equation}

\begin{lemma}[Majority expansion]\label{lem:majority}
Write
\begin{equation}\label{eq:generating}
F_\rho(v)=\sech^2\!\bigl(\rho\sqrt{1-v}\bigr)
=\sum_{k\ge0}d_kv^k.
\end{equation}
The coefficients are nonnegative and sum to one. Moreover
\[
\pi_k=\frac{\rho}{\tanh\rho}\frac{d_k}{C_k}
\]
is a probability distribution and
\begin{equation}\label{eq:majorityexpansion}
G_\rho(z)=\sum_{k\ge0}\pi_kM_k(z),\qquad -1\le z\le1.
\end{equation}
\end{lemma}
\paragraph{Proof idea.}
The cosine product gives nonnegative coefficients for a generating function.
Integrating the derivatives of odd-majority polynomials recovers tanh;
evaluation at the endpoint proves that the mixture weights sum to one.

\begin{proof}
The cosh product \cite{dlmf} gives
\[
F_\rho(v)=\prod_{j\ge0}
 \left(\frac{1-\eta_j}{1-\eta_jv}\right)^2,\qquad
\eta_j=\frac{\lambda_j}{1+\lambda_j},\quad
\lambda_j=\frac{4\rho^2}{\pi^2(2j+1)^2}.
\]
Each factor has nonnegative coefficients. The product converges
locally uniformly before its first pole, and
\(\sum_j\eta_j/(1-\eta_j)=\sum_j\lambda_j<\infty\).
It is a probability generating function with \(F_\rho(1)=1\).
For \(z\in[0,1]\), Tonelli's theorem and
\eqref{eq:majderivative} give
\[
\sum_k\pi_kM_k(z)
=\int_0^z\frac{\rho}{\tanh\rho}F_\rho(1-u^2)\,du
=G_\rho(z).
\]
At \(z=1\), this proves \(\sum_k\pi_k=1\).
Oddness gives negative \(z\).
\end{proof}

Sample \(K\) from \(\pi\), then request signs until one sign has
occurred \(K+1\) times. This gives the full majority's output while
sometimes using fewer than \(2K+1\) signs.

\begin{lemma}[Sequential majority cost]\label{lem:arity}
Its expected source count satisfies
\begin{equation}\label{eq:sequentialarity}
A_{\rm seq}(\rho,z)\le
1+\left(1-\frac{z^2}{3}\right)\rho^2+\rho^4.
\end{equation}
\end{lemma}
\paragraph{Proof idea.}
For each odd majority, sum the probabilities that a further vote is needed.
An integral identity expresses this cost in the same nonnegative
coefficients used by the factory and gives a uniform radius bound.

\begin{proof}
The identity \(C_k^{-1}=\int_0^1(1-v^2)^k\,dv\) follows by
integration by parts. The expected full-majority cost is therefore
\[
A(\rho)=\frac{\rho}{q}\int_0^1
\left[F_\rho(1-v^2)+2(1-v^2)F_\rho'(1-v^2)\right]\,dv.
\]
Put \(g(v)=\sech^2(\rho v)\).
Integration by parts applied to \(g(v)-1=O(v^2)\) gives
\begin{equation}\label{eq:arityintegral}
A(\rho)=\frac{\rho}{q}\left(1+
\int_0^1\frac{\tanh^2(\rho v)}{v^2}\,dv\right)
\le\frac{\rho}{q}(1+\rho^2).
\end{equation}
The singularity at zero is removable.
For \(K=1\), sequential majority costs \(5/2-z^2/2\), saving
\((1+z^2)/2\) source requests. Also
\(\pi_1=(2/3)\rho^2\sech^2\rho\).
Saving only this contribution gives
\[
A_{\rm seq}\le\frac{\rho}{q}(1+\rho^2)
-\frac{1+z^2}{3}\rho^2\sech^2\rho.
\]
Use \(\sech^2\rho\ge1-\rho^2\) and
\begin{equation}\label{eq:xcothx}
u\coth u\le1+\frac{u^2}3,\qquad u>0.
\end{equation}
For the latter, \((3+u^2)\sinh u-3u\cosh u\) vanishes at zero
and has derivative \(u(u\cosh u-\sinh u)\ge0\).
Substitution proves
\[
A_{\rm seq}\le
1+\left(1-\frac{z^2}{3}\right)\rho^2
+\frac{2+z^2}{3}\rho^4,
\]
which implies the claim.
\end{proof}

\subsection{Biases and exact stopping}\label{sec:bias}

Define \(T_\theta(u)=(u+\theta)/(1+\theta u)\).
The addition formula implies
\begin{equation}\label{eq:biasidentity}
\tanh(a+\rho z)=m+R\,T_{tq}(G_\rho(z)).
\end{equation}
For \(\theta\ge0\), put
\(\lambda=(1-\theta)/(1+\theta)\).
Request a source sign. Return on \(+1\); on \(-1\), return that
sign with probability \(\lambda\), otherwise repeat.
If the source's plus probability is \(p\), a trial succeeds with
\(\sigma=p+(1-p)\lambda\).
The accepted plus probability is \(p/\sigma\), with signed mean
\(T_\theta(2p-1)\).
For negative \(\theta\), complement the source and the output.
Its sign is the sign of the stored rational \(a\).

\begin{lemma}[Work of reached trials]\label{lem:stopping}
Suppose reaching trial \(j\) is determined before its fresh
randomness. If its conditional expected nonnegative cost is at
most \(B\), then
\begin{equation}\label{eq:stoppedcost}
\E\sum_{j=1}^T C_j\le B\sum_{j\ge1}\PP(T\ge j)=B\,\E T.
\end{equation}
For independent trials of success probability \(\sigma>0\),
\(\E T=1/\sigma\).
\end{lemma}
\paragraph{Proof idea.}
A trial is reached before its own randomness is exposed. Its expected work
can therefore be charged using the probability of reaching it, even when its
duration and its returned value are correlated.

\begin{proof}
Condition on the history before each trial, sum over finite
prefixes, and use monotone convergence. In the independent case
\(\PP(T\ge j)=(1-\sigma)^{j-1}\).
A trial's cost may be correlated with its outcome.
\end{proof}

\begin{lemma}[Joint majority and race cost]\label{lem:offspring}
A normalized call has expected children at most
\begin{equation}\label{eq:offspring}
\frac{\tanh\rho}{r_\ell}
\exp\!\left(\frac53\rho^2+\rho^4\right).
\end{equation}
A uniform layer bound is
\(B_\ell=\exp(\frac53(gr_{\ell-1})^2+(gr_{\ell-1})^4)\).
The root's expected children are at most \(r_DB_D\).
\end{lemma}
\paragraph{Proof idea.}
Multiply the mean cost of one reached majority by the mean number of race
trials. The acceptance probability cancels the bias denominator, leaving a
bound that depends on the row radius.

\begin{proof}
Assume \(t\ge0\), put \(u=\tanh(\rho z)\), and combine the gate,
the exact race mean, and Lemma~\ref{lem:stopping}:
\begin{equation}\label{eq:jointcost}
\E[\hbox{children}]
=\frac q{r_\ell} A_{\rm seq}(\rho,z)
\frac{1-t^2}{(1-tq)(1+tu)}.
\end{equation}
The last factor is
\[
K=\frac{\cosh\rho\,\cosh(\rho z)}
{\cosh(a-\rho)\cosh(a+\rho z)}
\le\frac{\cosh\rho\,\cosh(\rho z)}
{\cosh^2(\rho(1+z)/2)}.
\]
The denominator is minimized at \(a=\rho(1-z)/2\), by the
product-to-sum formula.
Since \((\log\cosh)''=\sech^2\le1\), its midpoint
second-difference bound yields
\(\log K\le\rho^2(1-z)^2/4\).
For negative \(t\), replace \(z\) by \(-z\). The arity is even
in \(z\), hence
\[
\log(A_{\rm seq}K)\le
\left(1-\frac{z^2}{3}+\frac{(1+|z|)^2}{4}\right)\rho^2+\rho^4
\le\frac53\rho^2+\rho^4.
\]
The prefactor in \eqref{eq:offspring} is at most one by
\eqref{eq:rounded}.
At the root the gate is \(R\), retaining the factor \(r_D\).
\end{proof}

Almost sure termination follows by induction on the layer.
The majority index is finite almost surely, lower-layer requests
terminate almost surely, and the race has positive success
probability. Known scalar coins terminate almost surely by
Section~\ref{sec:bits}. This reasoning precedes the expected-work
bound, so finiteness of the expectation is not assumed.
